\documentclass[11pt]{article}
\usepackage[margin=1in]{geometry}
\usepackage{amsmath,amssymb,amsthm,mathtools}
\usepackage{bm}
\usepackage{xcolor}
\usepackage{enumitem}

\begin{document}

\section*{Par 1: Uniform are Linear}

We will show that for every \(S\subseteq \Omega\),
\[
\{v_i:i\in S\}\text{ is linearly independent}
\quad \Longleftrightarrow \quad
|S|\le k.
\]
This proves that the uniform matroid is linear.


Let \(S=\{i_1,\dots,i_t\}\subseteq \Omega\), where \(t=|S|\).

We prove both directions.

\medskip
(1) If \(t>k\), then \(\{v_i:i\in S\}\) is linearly dependent.

Since each \(v_i\in \mathbb R^k\), the set \(\{v_{i_1},\dots,v_{i_t}\}\) consists of \(t\) vectors in the \(k\)-dimensional space \(\mathbb R^k\). If \(t>k\), then no set of \(t\) vectors in \(\mathbb R^k\) can be linearly independent. Hence \(\{v_i:i\in S\}\) is linearly dependent.

(2) If \(t\le k\), then \(\{v_i:i\in S\}\) is linearly independent.

Consider the \(k\times t\) matrix:
\[
A=
\begin{bmatrix}
1 & 1 & \cdots & 1\\
x_{i_1} & x_{i_2} & \cdots & x_{i_t}\\
x_{i_1}^2 & x_{i_2}^2 & \cdots & x_{i_t}^2\\
\vdots & \vdots & \ddots & \vdots\\
x_{i_1}^{k-1} & x_{i_2}^{k-1} & \cdots & x_{i_t}^{k-1}
\end{bmatrix}.
\]
We could take the first \(t\) rows of \(A\) which gives the \(t\times t\) submatrix
\[
B=
\begin{bmatrix}
1 & 1 & \cdots & 1\\
x_{i_1} & x_{i_2} & \cdots & x_{i_t}\\
x_{i_1}^2 & x_{i_2}^2 & \cdots & x_{i_t}^2\\
\vdots & \vdots & \ddots & \vdots\\
x_{i_1}^{t-1} & x_{i_2}^{t-1} & \cdots & x_{i_t}^{t-1}
\end{bmatrix}.
\]
Its transpose is the Vandermonde matrix
\[
B^{T}=
\begin{bmatrix}
1 & x_{i_1} & x_{i_1}^2 & \cdots & x_{i_1}^{t-1}\\
1 & x_{i_2} & x_{i_2}^2 & \cdots & x_{i_2}^{t-1}\\
\vdots & \vdots & \vdots & \ddots & \vdots\\
1 & x_{i_t} & x_{i_t}^2 & \cdots & x_{i_t}^{t-1}
\end{bmatrix}.
\]
By the Vandermonde determinant formula given in the hint,
\[
\det(B^T)=\prod_{1\le p<q\le t}(x_{i_q}-x_{i_p}).
\]
Since \(x_1,\dots,x_n\) are distinct, the numbers
\[
x_{i_1},\dots,x_{i_t}
\]
are also distinct, so every \(x_{i_q}-x_{i_p}\) is nonzero. Therefore
\[
\det(B^T)\neq 0.
\]
Hence \(\det(B)\neq 0\) as well, since \(\det(B)=\det(B^T)\),its columns are linearly independent. Thus \(B\) is invertible

Since \(B\) is a submatrix of \(A\), the columns of \(A\) are also linearly independent. Therefore
\[
\{v_{i_1},\dots,v_{i_t}\}
\]
is linearly independent.

\medskip
Combining the two parts, we have shown that
\[
\{v_i:i\in S\}\text{ is linearly independent}
\quad \Longleftrightarrow \quad
|S|\le k.
\]
Since \(S\in\mathcal I\) when \(|S|\le k\), it follows that
\[
S\in\mathcal I
\quad \Longleftrightarrow \quad
\{v_i:i\in S\}\text{ is linearly independent}.
\]
Therefore the uniform matroid is linear.



\section*{Part 2: Graphical are Linear}
We will show that for any $S\subseteq \Omega$,
\[
\{v_e:e\in S\}\text{ is linearly independent}
\quad\Longleftrightarrow\quad
S\text{ does not contain a cycle}.
\]

\noindent (1) If $S$ contains a cycle, then $\{v_e:e\in S\}$ is linearly dependent
\\
We first look at a cycle of size $3$, following the hint.

Consider the cycle on nodes $\{1,2,3\}$ with oriented edges
\[
e_1:1\to 2,\qquad e_2:1\to 3,\qquad e_3:3\to 2.
\]
Then
\[
v_{e_1}=(1,-1,0),\qquad
v_{e_2}=(1,0,-1),\qquad
v_{e_3}=(0,1,-1).
\]
Let
\[
A=[\,v_{e_1}\ \ v_{e_2}\ \ v_{e_3}\,]
=
\begin{bmatrix}
1&1&0\\
-1&0&1\\
0&-1&-1
\end{bmatrix}.
\]
To check whether these column vectors are linearly independent, solve
\[
x_1(1,-1,0)+x_2(1,0,-1)+x_3(0,1,-1)=(0,0,0).
\]
Equivalently,
\[
A
\begin{bmatrix}
x_1\\x_2\\x_3
\end{bmatrix}
=
\begin{bmatrix}
0\\0\\0
\end{bmatrix},
\]
that is,
\[
\begin{cases}
x_1+x_2=0,\\
-x_1+x_3=0,\\
-x_2-x_3=0.
\end{cases}
\]
Take for example
\[
x=
\begin{bmatrix}
1\\-1\\1
\end{bmatrix}\neq
\begin{bmatrix}
0\\0\\0
\end{bmatrix}.
\]
Then
\[
A x=0.
\]
So the homogeneous system has a nontrivial solution. Hence the column vectors of $A$ are linearly dependent. Equivalently,
\[
\det(A)=0,
\]
so $A$ is not invertible and its columns are linearly dependent.

Therefore, a cycle of size $3$ gives linear dependence.

\vspace{0.5em}

Now consider a general cycle
\[
u_1,u_2,\dots,u_{k},u_1
\]
with edges
\[
e_1=(u_1,u_2),\quad
e_2=(u_2,u_3),\quad
\dots,\quad
e_{k-1}=(u_{k-1},u_k),\quad
e_k=(u_k,u_1).
\]
Orient these edges consistently around the cycle. Then each vertex appears exactly once with coefficient $+1$ and exactly once with coefficient $-1$, so
\[
v_{e_1}+v_{e_2}+\cdots+v_{e_k}=0.
\]
This is a nontrivial linear relation, so the set
\[
\{v_{e_1},v_{e_2},\dots,v_{e_k}\}
\]
is linearly dependent.

Hence, if $S$ contains a cycle, then $\{v_e:e\in S\}$ is linearly dependent.
\vspace{1em}
\\

\noindent (2) If $S$ does not contain a cycle, then $\{v_e:e\in S\}$ is linearly independent

Assume now that $S$ is acyclic. The subgraph formed by $S$ is a forest.  
We first prove the statement for a tree, by induction, following the hint.

\subsubsection*{Step 1: Trees}

We prove that if $S$ is the edge set of a tree, then $\{v_e:e\in S\}$ is linearly independent.

1.Base case.
Consider the tree with two nodes and one edge:
\[
1-2.
\]
Let
\[
e_1=(1,2).
\]
Then
\[
v_{e_1}=(1,-1).
\]
If
\[
x_1v_{e_1}=0,
\]
then
\[
x_1(1,-1)=(0,0),
\]
which implies $x_1=0$. So $\{v_{e_1}\}$ is linearly independent.

Using the same style as above, one may also write
\[
A=[\,v_{e_1}\,]
=
\begin{bmatrix}
1\\
-1
\end{bmatrix},
\qquad
A[x_1]=0
\quad\Longrightarrow\quad
x_1=0.
\]
So the base case is proved.

\paragraph{Small example.}
Take the tree
\[
1-2-3,
\]
with edges
\[
e_1=(1,2),\qquad e_2=(2,3).
\]
Then
\[
v_{e_1}=(1,-1,0),\qquad
v_{e_2}=(0,1,-1).
\]
So
\[
A=[\,v_{e_1}\ \ v_{e_2}\,]
=
\begin{bmatrix}
1&0\\
-1&1\\
0&-1
\end{bmatrix}.
\]
If
\[
A
\begin{bmatrix}
x_1\\x_2
\end{bmatrix}
=0,
\]
then
\[
x_1v_{e_1}+x_2v_{e_2}=0,
\]
that is,
\[
x_1(1,-1,0)+x_2(0,1,-1)=(0,0,0).
\]
Looking at the first coordinate gives $x_1=0$, and then the third coordinate gives $x_2=0$. Hence these columns are linearly independent.

\paragraph{Inductive step.}
Let $S$ be the edge set of a tree with $n$ nodes.

Since every finite tree has a leaf, there exists a leaf $u$ in the tree. Let $e$ be the unique edge incident to $u$.

Suppose
\[
\sum_{i\in S} x_i v_i=0.
\]
We only consider the coordinate corresponding to the leaf $u$. Since $u$ is a leaf, only the vector corresponding to the unique incident edge $e$ can be nonzero in that coordinate. Thus, in the $u$-th coordinate,
\[
x_e\, v_e(u)=0.
\]
But $v_e(u)=1$ or $-1$, hence $v_e(u)\neq 0$, so necessarily
\[
x_e=0.
\]

Therefore, the coefficient of the edge $e$ must be zero. We can delete this edge and the leaf $u$. The remaining graph is still a tree, now with $n-1$ nodes. By the induction hypothesis, all remaining coefficients are also zero.

Hence every coefficient in
\[
\sum_{i\in S} x_i v_i=0
\]
must be zero. Therefore, the vectors $\{v_e:e\in S\}$ are linearly independent for any tree.

\subsubsection*{Step 2: Extend from trees to forests}

Now let $S$ be a forest. Then $S$ is a disjoint union of trees:
\[
S=S_1\cup S_2\cup\cdots\cup S_k,
\]
where each $S_j$ is the edge set of a tree.

Suppose
\[
\sum_{e\in S} x_e v_e=0.
\]
We can write this as
\[
\sum_{j=1}^k\ \sum_{e\in S_j} x_e v_e=0.
\]
Fix one component, say $S_j$. Since the forest is a disjoint union, the vectors outside of $S_j$ are zero on the coordinates which have the vertices as $j$-th tree. Therefore, if we only look at the coordinates of the $j$-th tree, we get
\[
\sum_{e\in S_j} x_e v_e=0.
\]
But $S_j$ is a tree, and from Step 1 we have proved that the corresponding vectors are linearly independent. Hence
\[
x_e=0\qquad \text{for all } e\in S_j.
\]
Since $j$ was arbitrary, this is true for every component. Therefore all coefficients are zero, so $\{v_e:e\in S\}$ is linearly independent.
\\
We proved both directions:

\begin{itemize}[leftmargin=2em]
    \item if $S$ contains a cycle, then $\{v_e:e\in S\}$ is linearly dependent;
    \item if $S$ is acyclic, then $\{v_e:e\in S\}$ is linearly independent.
\end{itemize}
Hence
\[
\{v_e:e\in S\}\text{ is linearly independent}
\quad\Longleftrightarrow\quad
S\text{ does not contain a cycle}.
\]
Therefore the graphical matroid is linear.

\end{document}